% GENERATED FILE. Edit canonical lessons, then run npm run generate:books.
% canonical-source-hash: fnv1a64:a809383a5a17d891
% canonical-lessons: GU-C407-pardes, GU-C407-vatan, GU-C407-dolar, GU-C407-pagpala, GU-C407-hijravu

\chapter{Far from Home}
\label{ch:gu-far-from-home}
\hlchaptermodality{407}

\begin{chapteropening}
\textbf{By the end of this chapter:} \emph{I can talk about living abroad and missing home.}

Complete the last lesson of chapter 407: I can talk about living abroad and missing home.
\end{chapteropening}

\section[\texorpdfstring{\gu{પરદેશ} (pardeś)}{પરદેશ (pardeś)}]{\gu{પરદેશ} (pardeś) — abroad; a foreign country}

\label{lesson:GU-C407-pardes}



\begin{quote}
Before the new one: say the Gujarati for an official fee, then the Gujarati for to empty out.
\end{quote}

\subsection*{You'll want to know: \gu{પરદેશ}}
\textbf{\gu{પરદેશ}} — \emph{pardeś} — \textquotedblleft{}abroad; a foreign country\textquotedblright{}.

\begin{quote}
\gu{મારો} \gu{ભાઈ} \gu{ત્રણ} \gu{વરસથી} \gu{પરદેશમાં} \gu{રહે} \gu{છે}.
\end{quote}

\begin{quote}
\emph{māro bhāī traṇ varasthī pardeśmā̃ rahe che.}
\end{quote}

\begin{quote}
My brother has lived abroad for three years.
\end{quote}

A masculine noun.

\begin{cousinweb}
\textbf{\gu{પરદેશ}} (\emph{pardeś}) comes from Sanskrit \emph{paradeśa}, \textquotedblleft{}another or a foreign country\textquotedblright{} (Monier-Williams): any land that is not your own.
\end{cousinweb}

\begin{grammarlens}[title={What you've built}]
One more word for getting around.
\end{grammarlens}

\subsection*{Your turn}
\begin{itemize}
\raggedright
  \item \emph{Say it:} \emph{pardeś}
  \item \emph{Say it:} \emph{pardeś}, once more
  \item \emph{Say it:} \emph{māro bhāī traṇ varasthī pardeśmā̃ rahe che.} — My brother has lived abroad for three years.
  \item \emph{Recall:} say the Gujarati for an official fee, then the Gujarati for to empty out, then say \emph{pardeś} again
\end{itemize}

\begin{tcolorbox}[breakable,colback=teal!4,colframe=teal!35!black,title={Before you move on}]
What does \gu{પરદેશ} mean? (\textquotedblleft{}Abroad; a foreign country\textquotedblright{}.) Say it once more.
\end{tcolorbox}
\section[\texorpdfstring{\gu{વતન} (vatan)}{વતન (vatan)}]{\gu{વતન} (vatan) — a homeland; one's home town}

\label{lesson:GU-C407-vatan}



\begin{quote}
Before the new one: say the Gujarati for to empty out, then the Gujarati for abroad.
\end{quote}

\subsection*{You'll want to know: \gu{વતન}}
\textbf{\gu{વતન}} — \emph{vatan} — \textquotedblleft{}a homeland; one's home town\textquotedblright{}.

\begin{quote}
\gu{દિવાળીમાં} \gu{બધા} \gu{વતન} \gu{પાછા} \gu{આવે} \gu{છે}.
\end{quote}

\begin{quote}
\emph{divāḷīmā̃ badhā vatan pāchā āve che.}
\end{quote}

\begin{quote}
At Diwali everyone comes back to their home town.
\end{quote}

A neuter noun.

\begin{cousinweb}
\textbf{\gu{વતન}} (\emph{vatan}) is usually traced to Arabic \emph{waṭan}, \textquotedblleft{}a homeland, a dwelling place\textquotedblright{}, which came through Persian: the place you belong to, wherever you live now.
\end{cousinweb}

\begin{grammarlens}[title={What you've built}]
One more word for getting around.
\end{grammarlens}

\subsection*{Your turn}
\begin{itemize}
\raggedright
  \item \emph{Say it:} \emph{vatan}
  \item \emph{Say it:} \emph{vatan}, once more
  \item \emph{Say it:} \emph{divāḷīmā̃ badhā vatan pāchā āve che.} — At Diwali everyone comes back to their home town.
  \item \emph{Recall:} say the Gujarati for to empty out, then the Gujarati for abroad, then say \emph{vatan} again
\end{itemize}

\begin{tcolorbox}[breakable,colback=teal!4,colframe=teal!35!black,title={Before you move on}]
What does \gu{વતન} mean? (\textquotedblleft{}A homeland; one's home town\textquotedblright{}.) Say it once more.
\end{tcolorbox}
\section[\texorpdfstring{\gu{ડોલર} (ḍolar)}{ડોલર (ḍolar)}]{\gu{ડોલર} (ḍolar) — a dollar}

\label{lesson:GU-C407-dolar}



\begin{quote}
Before the new one: say the Gujarati for abroad, then the Gujarati for a homeland.
\end{quote}

\subsection*{You'll want to know: \gu{ડોલર}}
\textbf{\gu{ડોલર}} — \emph{ḍolar} — \textquotedblleft{}a dollar\textquotedblright{}.

\begin{quote}
\gu{એક} \gu{ડોલરના} \gu{કેટલા} \gu{રૂપિયા} \gu{મળે} \gu{છે}?
\end{quote}

\begin{quote}
\emph{ek ḍolarnā keṭlā rūpiyā maḷe che?}
\end{quote}

\begin{quote}
How many rupees do you get for a dollar?
\end{quote}

A masculine noun.

\begin{cousinweb}
\textbf{\gu{ડોલર}} (\emph{ḍolar}) is borrowed from English \emph{dollar}.
\end{cousinweb}

\begin{grammarlens}[title={What you've built}]
One more word for getting around.
\end{grammarlens}

\subsection*{Your turn}
\begin{itemize}
\raggedright
  \item \emph{Say it:} \emph{ḍolar}
  \item \emph{Say it:} \emph{ḍolar}, once more
  \item \emph{Say it:} \emph{ek ḍolarnā keṭlā rūpiyā maḷe che?} — How many rupees do you get for a dollar?
  \item \emph{Recall:} say the Gujarati for abroad, then the Gujarati for a homeland, then say \emph{ḍolar} again
\end{itemize}

\begin{tcolorbox}[breakable,colback=teal!4,colframe=teal!35!black,title={Before you move on}]
What does \gu{ડોલર} mean? (\textquotedblleft{}A dollar\textquotedblright{}.) Say it once more.
\end{tcolorbox}
\section[\texorpdfstring{\gu{પગપાળા} (pagpāḷā)}{પગપાળા (pagpāḷā)}]{\gu{પગપાળા} (pagpāḷā) — on foot}

\label{lesson:GU-C407-pagpala}



\begin{quote}
Before the new one: say the Gujarati for a homeland, then the Gujarati for a dollar.
\end{quote}

\subsection*{You'll want to know: \gu{પગપાળા}}
\textbf{\gu{પગપાળા}} — \emph{pagpāḷā} — \textquotedblleft{}on foot\textquotedblright{}.

\begin{quote}
\gu{અમે} \gu{સ્ટેશન} \gu{સુધી} \gu{પગપાળા} \gu{ગયા}.
\end{quote}

\begin{quote}
\emph{ame sṭeśan sudhī pagpāḷā gayā.}
\end{quote}

\begin{quote}
We went to the station on foot.
\end{quote}

\begin{cousinweb}
\textbf{\gu{પગપાળા}} (\emph{pagpāḷā}) is built from \gu{પગ}, a foot, and \emph{pāḷũ}, \textquotedblleft{}on foot\textquotedblright{}, which CDIAL lists for Gujarati under \emph{*pādala}, \textquotedblleft{}belonging to the feet\textquotedblright{}, a form it reconstructs (CDIAL 8066).
\end{cousinweb}

\begin{grammarlens}[title={What you've built}]
One more word for getting around.
\end{grammarlens}

\subsection*{Your turn}
\begin{itemize}
\raggedright
  \item \emph{Say it:} \emph{pagpāḷā}
  \item \emph{Say it:} \emph{pagpāḷā}, once more
  \item \emph{Say it:} \emph{ame sṭeśan sudhī pagpāḷā gayā.} — We went to the station on foot.
  \item \emph{Recall:} say the Gujarati for a homeland, then the Gujarati for a dollar, then say \emph{pagpāḷā} again
\end{itemize}

\begin{tcolorbox}[breakable,colback=teal!4,colframe=teal!35!black,title={Before you move on}]
What does \gu{પગપાળા} mean? (\textquotedblleft{}On foot\textquotedblright{}.) Say it once more.
\end{tcolorbox}
\section[\texorpdfstring{\gu{હિજરાવું} (hijrāvũ)}{હિજરાવું (hijrāvũ)}]{\gu{હિજરાવું} (hijrāvũ) — to pine for home; to grieve inwardly}

\label{lesson:GU-C407-hijravu}



\begin{quote}
Before the new one: say the Gujarati for a dollar, then the Gujarati for on foot.
\end{quote}

\subsection*{You'll want to know: \gu{હિજરાવું}}
\textbf{\gu{હિજરાવું}} — \emph{hijrāvũ} — \textquotedblleft{}to pine for home; to grieve inwardly\textquotedblright{}.

\begin{quote}
\gu{પરદેશમાં} \gu{દાદી} \gu{વતન} \gu{માટે} \gu{હિજરાતી} \gu{હતી}.
\end{quote}

\begin{quote}
\emph{pardeśmā̃ dādī vatan māṭe hijrātī hatī.}
\end{quote}

\begin{quote}
Abroad, grandmother used to pine for her home town.
\end{quote}

Past (she …): \gu{તે} \gu{ઘણી} \gu{હિજરાઈ}, \emph{te ghaṇī hijrāī}.

\begin{cousinweb}
\textbf{\gu{હિજરાવું}} (\emph{hijrāvũ}) is usually traced to Persian \emph{hijrān}, \textquotedblleft{}separation\textquotedblright{}, a word of Arabic origin: the ache of being far from people and places you love.
\end{cousinweb}

\begin{grammarlens}[title={What you've built}]
One more word for getting around.
\end{grammarlens}

\subsection*{Your turn}
\begin{itemize}
\raggedright
  \item \emph{Say it:} \emph{hijrāvũ}
  \item \emph{Say it:} \emph{hijrāvũ}, once more
  \item \emph{Say it:} \emph{pardeśmā̃ dādī vatan māṭe hijrātī hatī.} — Abroad, grandmother used to pine for her home town.
  \item \emph{Recall:} say the Gujarati for a dollar, then the Gujarati for on foot, then say \emph{hijrāvũ} again
\end{itemize}

\begin{tcolorbox}[breakable,colback=teal!4,colframe=teal!35!black,title={Before you move on}]
What does \gu{હિજરાવું} mean? (\textquotedblleft{}To pine for home; to grieve inwardly\textquotedblright{}.) Say it once more.
\end{tcolorbox}
